\subsection{Flag nefness outside the correction primes}
\label{sec:geometry-uncorrected-flag-nefness}

\begin{lemma}[Degrees from the Hodge and conjugate filtrations]
\label{lem:geometry-two-filtration-nefness}
On a smooth projective connected curve \(C\) over an algebraically
closed field of characteristic \(p>2\), suppose the vector bundles
\(L,H\) have exact sequences
\[
 0\to L^\vee\to H\to L\to0,\qquad
 0\to F_C^*L\to H\to F_C^*L^\vee\to0.
\]
Here \(F_C\) is absolute Frobenius. Suppose these sequences respect
the split unramified CM coefficient action: \(L=\bigoplus_\sigma
L_\sigma\), duality pairs conjugate labels, and Frobenius permutes the
labels by \(\varphi(\sigma)=\sigma^p\). Put \(e=-\deg L\). Then
\[
 e\ge0,\qquad \deg W\le e/(p-2)\quad(W\subset L\text{ a subbundle}),
 \qquad \deg L_\sigma=\deg L_{\bar\sigma}.
 \tag{UN1}
\]
\end{lemma}
\begin{proof}
Let \(P\) be the maximum degree of saturated subbundles of \(L\),
including the zero subbundle. It exists: dualizing a generating map
for a sufficiently positive twist of \(L^\vee\) embeds \(L\) in a
sum of positive line bundles, and exterior powers give upper bounds
for the integer degrees in each rank. In particular \(P\ge0\).

Pull a subbundle of degree \(P\) back by Frobenius and inject it in
\(H\) by the second sequence. Its degree is \(pP\). Project to \(L\)
by the first sequence. The image has degree at most \(P\), since
saturation can only increase degree. The kernel in \(L^\vee\) has
degree at most \(e+P\): for a saturated subbundle \(K\subset L^\vee\),
\(\deg K=-\deg L+\deg (L^\vee/K)^\vee\le e+P\), since the last
dual quotient embeds as a subbundle of \(L\). This bound also holds
without saturation. Additivity of degree gives \(pP\le e+2P\).

Write \(d_\sigma=\deg L_\sigma\) and \(h_\sigma=\deg H_\sigma\).
The first sequence gives \(h_\sigma=d_\sigma-d_{\bar\sigma}\),
and the second gives \(h_{\varphi(\sigma)}=p h_\sigma\).
The permutation has finite orbits, so every \(h_\sigma=0\).
This proves (UN1).
\end{proof}

\begin{proposition}[The actual uncorrected finite-place flag bound]
\label{prop:geometry-uncorrected-flag-nefness}
For the coordinated original integral models, at each prime
\(p\notin S\), \(p>2\), \(p\nmid\operatorname{disc}(E)\), the root
model satisfies \(\mathrm{NefAt}(p,\mu,0)\) for every dominant
weight satisfying the original threshold
\[
 (p-2)(\mu_{\tau,j}-\mu_{\tau,j+1})
 \ge 4 f\|\mu\|_\infty\qquad(\text{every CM type label }\tau).
 \tag{UN2}
\]
Here \(f=[E:\mathbb Q]/2\), the weights are indexed from \(1\) to
\(N\), and \(j,N-j\) are the positive ranks of the paired tangent
eigenspaces. This supplies \(M2\) for the actual good tower,
for which the prime hypotheses follow from
\(p\notin\mathrm{InS}(B)\).
\end{proposition}
\begin{proof}
We first supply the exact sequences on the original test curves,
including curves contained in the boundary. The original unimodular
symplectic lattice with its \(O_E\)-action has a hyperspecial unitary
similitude model at this prime: \(O_E\otimes\mathbb Z_p\) is finite
etale, and the perfect pairing and involution give the reductive
similitude group after unramified splitting. The underlying principally
polarized PEL datum is unchanged (the primitive biextension's diagonal
factor \(2\) is a unit here). Forget the \(p\)-primary part of the
principal level, if present. The remaining level is neat by
\(\mathrm{GoodAt}(p)\). The coordinated level transitions and canonical
normalization construction identify the actual semiabelian family
with the pullback from this prime-to-\(p\) PEL toroidal model, also on
the global dominating model and flag modifications by their
cartesian comparisons.

Use Goldring--Koskivirta, \emph{Strata Hasse invariants, Hecke
algebras and Galois representations}, Invent. Math. 217 (2019),
887--984, Theorem 6.2.1 and Lemma 6.2.3, with their canonical
de Rham and Hodge bundles (Section 6.1). Refine the rational cones
to admit their log-integral Siegel compactification. For the
hyperspecial PEL embedding the theorem gives the Hodge and conjugate
filtrations with their coefficient tensors and Frobenius graded
isomorphisms on the whole toroidal model. Its canonical Hodge bundle
is the invariant cotangent bundle of the semiabelian family.
In de Rham homology the Hodge quotient is \(L=\operatorname{Lie}(G)\)
and the kernel is \(L^\vee\) with conjugate coefficient action by
the polarization. Dualizing the cohomological conjugate filtration
gives subbundle \(F_C^*L\) and quotient \(F_C^*L^\vee\).
The Frobenius maps raise residue embeddings to their \(p\)-th
powers. These are the two sequences in
Lemma~\ref{lem:geometry-two-filtration-nefness}.

For precision, the auxiliary fan refinement does not change the
arithmetic model or its degree. For an integral projective curve
of the original special fibre, base change to the algebraic closure
of the residue field and normalize a component. Pull back the proper
surjective refinement along its map to the prime-to-\(p\) model.
A closed point in the nonempty generic fibre of that pullback has
finite residue degree over the curve function field. Its closure,
normalized, is a smooth projective curve finite and surjective over
the chosen curve. The family, coefficient action and Hodge bundle
pull back along this map, and line-bundle degrees multiply by the
positive degree of the map. It suffices to prove the sign on this
test curve. In particular the special-fibre curve need not have an
abelian generic point. No uniform bound on the auxiliary curve
degree is needed for nefness.

Since \(p\notin S\), the patched tangent lattices agree with the
original tangent eigenlattices on the entire curve. The integral
idempotents of the etale coefficient algebra split the tangent
bundle and preserve the given subbundle flags. The root line
numbered \(i\le j\) is the \(i\)-th graded line of \(L_\tau\);
for \(i>j\) it is the dual of graded line \(N+1-i\) of
\(L_{\bar\tau}\), precisely the original
\(\mathrm{rootIdx},\mathrm{rootSgn}\) prescription.

Write \(q=N-j\); let \(a_{\tau,i},b_{\tau,i}\) be the degrees of
the prefixes of these two flags, starting at zero. By (UN1) all
are at most \(P=e/(p-2)\), with
\(a_{\tau,j}=b_{\tau,q}=d_\tau\) and \(\sum_\tau d_\tau=-e/2\).
Let \(S_\tau\) be the sum of root-line degrees weighted by
\(\mu_{\tau,i}\). Abel summation in the two flags, suppressing
\(\tau\) on the weights, gives
\[
\begin{split}
 S_\tau={}&(\mu_j-\mu_{j+1})d_\tau
 +\sum_{i=1}^{j-1}(\mu_i-\mu_{i+1})a_{\tau,i}\\
 &+\sum_{i=j+1}^{N-1}(\mu_i-\mu_{i+1})b_{\tau,N-i}.
\end{split}
\]
All displayed differences are nonnegative. Put
\(g=\min_\tau(\mu_{\tau,j}-\mu_{\tau,j+1})\) and
\(M=\|\mu\|_\infty\). Because \(d_\tau\le P\),
\[
 S_\tau\le g d_\tau+2MP,\qquad
 \sum_\tau S_\tau\le -ge/2+2fMP\le0
\]
by (UN2). The original numerical class has coefficients
\(-\mu_{\tau,i}\) at the root lines and zero at the extra line
when its last parameter is \(0\). Its curve degree is therefore
\(-\sum_\tau S_\tau\ge0\), exactly the required nefness.
\end{proof}

Together with the corrected-flag theorem at primes in \(S\), this
establishes both finite-place flag-positivity suppliers. It does not
establish the two-level arithmetic Chern-class comparison, which
requires the exact vertical correction sequences and their Frobenius
relation rather than only numerical degrees on curves.
